% NOETHER PAPER 6 — KOREAN TRANSLATION-PRODUCER DRAFT — UNCHECKED
% Unit: T03-U21; current-authority whole lines 4777--4793
% Source slice: 938 LF UTF-8 bytes; SHA-256 F7F81CBB7440996B2B0390DA7388EEB590BB619D4B24498AB6DFD293BB48E214
% Pointer: NOETH-DE-AUTH-v007-20260804; SHA-256 A6A8FC8E5AC24ACAF49DFD55B4B58FA3DA882EF8C3FDD4D136220C8751045156
% German authority: NOETH-DE-ED-0001; 2,153,565 raw bytes; SHA-256 D1F06B311F6CBD991DD247D745DD9A72DDE326A20396DF43CFE0C8EDB1593CDB
% Producer translation only: no source/Korean/formula review, compilation, or rendering.
% Terminology debt: sukzessive Substitution→잇따른 대입, reduzierte Darstellung→기약표현, Rang erhalten→계수(階數)가 보존되다 are provisional.
% Cross-unit topology: this unit opens the \srcfn footnote; U22 continues and closes it. Concatenate U21 then U22 without inserted prose.
% Handoff state: UNCHECKED; all Korean wording and preserved TeX require an independent Korean checker.

\emph{수들}
\[
 a_n,\ a_{n-1},\ldots,\ a_{\rho+1}
\]
\emph{은 다음 대입에 의해}
\[
 x_n=a_n,\quad x_{n-1}=a_{n-1},\ldots,\quad x_{\rho+1}=a_{\rho+1}
\]
\emph{곱 $G(x)\cdot\Gamma(x)$가 영이 되지 않도록 --- 즉 계 (1)의 대수적 계수(階數)가 보존되도록 --- 택한다고 하자. 잇따른 대입들}
\[
\begin{array}{c}
[h(x)]_{x_n=a_n}=h^{(n-1)}(x);\quad
[h^{(n-1)}(x)]_{x_{n-1}=a_{n-1}}=h^{(n-2)}(x);\quad\cdots;\\[0.25em]
[h^{(\rho+1)}(x)]_{x_{\rho+1}=a_{\rho+1}}=h^*(x_1,x_2,\ldots,x_\rho)
\end{array}
\]
\emph{에서는 함수들 $h(x),h^{(n-1)}(x),\ldots,h^{(\rho+1)}(x)$을 각각 기약표현으로 만든다고 하자. 이 가정들 아래에서는 함수 $h^*(x_1\cdots x_\rho)$의 분자도 분모도 항등적으로 영이 될 수 없으며, 이 함수는 $\frac00$이 될 수도 없다.}\srcfn{**)}{잇따른 대입은 예컨대 앞 주석의 함수 $h(x)$에서 다음을 준다:
